
\section{Model checking based control in local space}
\label{ModelChackingLocal}


% \section{Evaluation}
\subsection{Overview}
% Input: Initial position, target position, set of obstacles
% Robot structure (xacro) -> Generate NuSMV -> check model, get counterexample -> 
% derive path -> compare with shortest path algorithm (implement Deikstra/ A*)
%
% Evaluation: comparison, 
% Validation: assembly simulation in Gazebo
A basic verification step for our model consists in checking whether a given point in 
space is reachable by the robot’s end-effector when the blocks are arranged 
in a specific initial configuration. 
Establishing such reachability opens the way for identifying potential stages 
of the assembly process, as illustrated in Figure \ref{fig:LocalAssemblyWithRMCandCE}.

\begin{figure}[h!]
    \centering
    \includegraphics[height=15cm]{LocalAssemblyWithRMCandCE.png} \caption{Self-replication process with potential model checking stages}
    \label{fig:LocalAssemblyWithRMCandCE}
\end{figure}

In this context, model checking can be employed to confirm that the 
robot is capable of executing fundamental operations 
required for self-replication: 
\begin{itemize}
    \item moving the end-effector from its current location to the position of a target part, 
    \item transporting the part to the designated assembly site, 
    \item moving the newly attached end-effector to connect with the next required block. 
\end{itemize}

In this chapter, we present a NuSMV model that encodes a chosen robot configuration, 
together with temporal logic specifications for verifying point and 
region reachability. Furthermore, we introduce Python-based tooling that 
derives counterexamples from the model checker and interprets them as 
executable control strategies for guiding the robot through the assembly process.


\subsection{Robot NuSMV model}
\label{ch:robotNuSMVmodel}

The robot structure as it shown in Figure \ref{fig:RobotStructureGazebo}
is represented as a modular composition of three distinct block types: a fixed 
base block, blocks with a yaw degree of freedom, and blocks with a pitch degree of 
freedom. Figure \ref{fig:RobotConfiguration} illustrates the hierarchical arrangement of 
these blocks and the flow of coordinate and orientation 
data between them. This modular abstraction allows the 
robot structure to be described in a uniform and 
scalable manner, independent of the number of blocks.
\begin{figure}[h!]
    \centering
    \includegraphics[width=16cm]{Robot NuSMV model.png} 
    \caption{Robot configuration NuSMV model}
    \label{fig:RobotConfiguration}
\end{figure}
% \textcolor{red}{TODO: 2. Correct diagram to SYSML block connection diagram. remove transition system}

As the lowest level of a system, as it shown in Figure \ref{fig:RobotConfiguration} we define a block.
Each subsequent block introduces a transformation of the block end position according to its type.
Each block has a certain type, in this work i consider three types of blocks: \\ 
The $block\_base$ defines the 
fixed starting point of the robot. It receives as input the global origin coordinates 
$ (x\_orig,y\_orig,z\_orig) $
and the initial orientation parameters 
$ (yaw\_orig,pitch\_orig) $
Its role is to provide a reference frame for 
all subsequent blocks. The base block outputs the coordinates 
$ (x\_end,y\_end,z\_end) $
 and orientation values that are propagated forward to the next block in the sequence.


Block types block\_yaw realise rotation around local axis Z, that oriented along the parent block
The $block\_yaw$ represents a joint capable of rotating around the vertical axis. 
Internally, the block maintains a discrete variable 
$ \phi $, which denotes the yaw angle. 

As illustrated in the Figure \ref{fig:BlockTS}, 
$\phi$ can evolve in steps of +1 or -1, modulo 36, corresponding to 10-degree discretization 
over the full circle of 360°. Thus, the block admits cyclic transitions that allow 
continuous rotation in both clockwise and counterclockwise directions. 
The block takes as input the parent coordinates 
$ (x\_par,y\_par,z\_par) $ and orientation $ (yaw\_par,pitch\_par) $,
applies the yaw rotation, and produces updated endpoint coordinates.
\begin{equation}
    \begin{cases}
        \phi \rightarrow \phi\\
        \phi \rightarrow (\phi+1)\ mod\ 36\\
        \phi\rightarrow(\phi-1)\ mod\ 36\\
    \end{cases}
\end{equation}

Block types block\_pitch realise rotation around local axis Y, that is perpendicular to Z often
the parent block.
Similarly, the $block\_pitch$ encodes a joint that rotates around the lateral axis 
of the preceding link. Its discrete state variable $\theta$:
\begin{equation}
    \begin{cases}
        \theta \rightarrow \theta\\
        \theta \rightarrow (\theta+1)\ mod\ 36\\
        \theta\rightarrow(\theta-1)\ mod\ 36 \\
    \end{cases} 
\end{equation}
In combination, these transitions capture the continuous ability to rotate up 
or down by discrete steps. 

\begin{figure}[h!]
    \centering
    \includegraphics[width=15cm]{Block transition system.png} 
    \caption{Robot block transition system for angle $\phi$}
    \label{fig:BlockTS}
\end{figure}

At the system level, the robot is built as a chain of such blocks. After the 
base block, one may attach any sequence of yaw and pitch blocks in arbitrary order, depending 
on the desired manipulator structure. The right-hand side of 
Figure \ref{fig:RobotConfiguration} shows an example with N chained blocks, where Block
0 connects directly to the base and Block N defines the robot’s end-effector. 
The outputs of the last block are new coordinates 
$ (x\_end,y\_end,z\_end) $ and updated orientation values.
Each block is parameterized by the output of its predecessor, creating a cascading 
dependency structure.

Thus, on the NuSMV model description level we have:
\begin{itemize}
    \item main block module - that describes a robot structure as a sequence of 
        block types as variables in VAR section.
    \item block\_base module - describing initial block of a sequence and locating on a ground plane.
    \item block\_yaw module - describing rotation around the parent Z axis (along the block)
        and changing the yaw angle, the NuSMV code for this module is shown in 
        Listing \ref{lst:NuSMVyawblock}.
    \item block\_pitch module - describing rotation around the parent Y axis (perpendicular to 
        the block) and changing the pitch angle.
\end{itemize}

\lstinputlisting[language=NuSMV, basicstyle=\ttfamily\scriptsize, 
numbers=left, firstline=57, lastline=90, 
caption={NuSMV code for block\_yaw module}, label={lst:NuSMVyawblock}]{code/robot_structure.smv}

MuSMV modules block\_base, block\_pitch sin, and cos NuSMV modules shown in Appendix A.
This design supports a natural inductive description of 
the robot in NuSMV: the same block type can be instantiated repeatedly, 
with only the input-output connections altered.

The modularity of the design has two important advantages.\\
First, it allows scalability, since additional joints can be incorporated 
simply by adding new instances of yaw or pitch modules. \\
Second, it supports module based model checking, because each block is described 
as a transition system with a finite number of states and transitions. 
The global state space of the robot is then obtained by synchronous product of 
the component transition systems, which can be directly 
analyzed using LTL or CTL specifications in NuSMV.

The dynamics of each yaw or pitch block can be represented as a finite-state transition system. 
Figure \ref{fig:BlockTS} illustrates this equivalence explicitly. 
On the left-hand side, the system is shown as a ring of discrete states 
corresponding to the possible angle values 
\[ \phi=0^\circ,10^\circ,…,350^\circ \]
From each state, there are transitions to the neighboring states at 
$ \phi+10^\circ $
 and 
 $ \phi-10^\circ $
thus modeling incremental rotations in either direction. Since the variable is defined modulo 36, the state 
$ \phi=350^\circ $
 transitions back to 
 $ \phi=0^\circ $
, completing the cycle. The result is a finite cyclic transition system with 36 nodes, 
which represents the discretized continuous rotation of the block.

On the right-hand side of Figure \ref{fig:BlockTS}, this same behavior is captured more compactly.
Here, the internal state variable $\phi$ is updated according to transition rules:

\[ \phi=\phi \quad \lor \quad \phi=(\phi+1)\ mod\ 36\ \lor \phi=(\phi-1)\ mod\ 36 \]

This concise representation preserves the semantics of the explicit ring but makes the 
module description more compact and scalable. When combined with coordinate update 
formulas, it fully captures the contribution of the block to the robot’s kinematics.

This transitions and angle changing for each block reflected by the following part  
of assign block, shown in Listing \ref{lst:AssignPart}.

\begin{lstlisting}[language=NuSMV, caption={Assign part of code 
    for block transitions}, label={lst:AssignPart}, float]
ASSIGN
    ...
    init(yaw)   := 0;
    next(yaw)   := { 
                    yaw,                
                    (yaw+1) mod 36,    
                    (yaw+35) mod 36    
                    };
\end{lstlisting}
For each block in a sequence of robot structure the initial value is assigned, and next state 
can be chosen nondeterministically from three possible transitions.

The full NuSMV model for modules: block\_base, block\_yaw, block\_pitch and main presented in Appendix A.

\subsection{Reachability specifications}
\label{ch:reachabilitySpec}
\subsubsection{Point reachability checking specifications:}

In order to verify whether the robot is capable of reaching a desired target point 
with its free end-effector, in three-dimensional space, we formalize the reachability 
check directly within the NuSMV model. 

In the DEFINE section we provide target point, end-effector coordinates variable 
and define error threeshold, this code fragment is shown in Listing \ref{lst:CheckPreparation}.

\begin{lstlisting}[language=NuSMV, 
    caption={End effector and target point coincidence condition}, 
    label={lst:CheckPreparation}]
DEFINE
...
error  := 1;
checkX := 10;
checkY := 15;
checkZ := 25;
endX_inside_limits :=(endX <=(checkX+error) & endX >=(checkX-error));
endY_inside_limits :=(endY <=(checkY+error) & endY >=(checkY-error));
endZ_inside_limits :=(endZ <=(checkZ+error) & endZ >=(checkZ-error));
\end{lstlisting}

Here, the variables endX, endY, and endZ represent the Cartesian coordinates 
of the robot end-effector, which are computed elsewhere in the model 
based on the joint angles and link lengths. 

The parameters checkX, checkY, and checkZ 
denote the coordinates of a target point that we wish to test for reachability. 

Since we can use only countable number of states and have to represent angles as  integers,
each transformation of angles to linear coordinates is not precise. 
Correspondingly the distance between actual and target points is between zero and linear 
delta calculated based on descretization value for angle.
For each spatial dimension, Boolean formulas $endX\_inside\_limits$, $endY\_inside\_limits$, 
and $endZ\_inside\_limits$ check whether the end-effector coordinate lies within a 
small interval around the corresponding target value, effectively 
defining a cube of side length 2, centered at the target.

The formal property is expressed in Linear Temoporal Logic (LTL) in Listing \ref{lst:LTLSpec}.
states that on every execution path, it will eventually (F) be 
the case that the robot’s end-effector is  within the specified tolerance region around 
the target point. In other words, the property enforces that the end-effector  
eventually can reach the target zone.

\begin{lstlisting}[language=NuSMV, caption={Reachability condition in LTL specification}, 
    label={lst:LTLSpec}]
    LTLSPEC 
    F (endX_inside_limits & endY_inside_limits & endZ_inside_limits)
\end{lstlisting}
The formal property is expressed in Computation Tree Logic (CTL) in Listing \ref{lst:CTLSpec}.

\begin{lstlisting}[language=NuSMV, caption={Reachability condition in CTL specification}, 
    label={lst:CTLSpec}]
CTLSPEC 
EF (endX_inside_limits & endY_inside_limits & endZ_inside_limits)
\end{lstlisting}

In this specification $EF$ means "there exists a path along 
which eventually holds following condition". 


\subsubsection{Region reachability checking specifications:}

We can check reachability of some 3D region, in our case a box by defining the 
variables checkX, checkY and checkZ as FROZENVAR. 
FROZENVAR section in NuSMV declares parameters that are chosen 
nondeterministically at starting time 
and then never change along a run for an entire execution (trace).
The Listing \ref{lst:RegionDefinition}  represents an example 
of a region: \\ $x=[-10,10],y=[-10,10],z=[0,20]$.

\begin{lstlisting}[language=NuSMV, caption={Region definition for a reachability checking}, 
    label={lst:RegionDefinition}]
FROZENVAR
    checkX : -10..10;
    checkY : -10..10;
    checkZ : 0..20;
\end{lstlisting}

Changing in the specification for the region checking using CTL logic 
shown in Listing \ref{lst:CTLSpec} is not necessary.
With the frozen variable region, the formula states: from every initial state 
(for every admissible target chosen at t=0), there exists a sequence 
of moves that eventually places the end-effector inside the region.

Because NuSMV checks CTL properties over all initial states
, this specification proves reachability for every target in the declared 
ranges (universal over targets, existential over paths). 

These specifications therefore provides a rigorous, formally verifiable 
mechanism to test spatial reachability of robotic manipulators in 
discrete state-space models. It can be 
generalized by altering the checkX, checkY, and 
checkZ parameters, and by varying the tolerance error,
to systematically analyze the reachable workspace of the system.

\subsection{Counterexample processing}
\label{ch:Counterexample}

Based on specifications described in previous section we use 
the modeling pattern in NuSMV for reachability: “eventually reach point P”, 
and negate that requirement to force the model checker to produce a counterexample 
that reaches target point P. 

In our self-replication robotic system this pattern can be implementated to obtain 
plausible traces on following stages of the self-replication process:
\begin{enumerate}
    \item To define the trajectory from current point to new part.
    \item To define the trajectory of transportation of the part to an assembly point.
    \item To define the trajectory for new robot to connect to the last part.
\end{enumerate}

These stages shown on the Figure \ref{fig:LocalAssemblyWithRM}, where input parameters 
shown in orange squares and counterexamples depicted with red arrow.

\begin{figure}[h!]
    \centering
    \includegraphics[height=16cm]{LocalAssemblyWithRMCandCE.png} 
    \caption{Local assembly process without model checking}
    \label{fig:LocalAssemblyWithRM}
\end{figure}

The formal property is expressed in Linear Temoporal Logic (LTL) in Listing \ref{lst:LTLSpecNeg}.
states that on every execution path, it will always (G) be 
the case that the robot’s end-effector is not within the specified tolerance region around 
the target point. In other words, the property enforces that the end-effector infinitely often 
leaves the target zone, meaning the robot cannot remain in that region permanently.

\begin{lstlisting}[language=NuSMV, caption={Reachability condition in LTL specification}, 
    label={lst:LTLSpecNeg}]
LTLSPEC 
G !(endX_inside_limits & endY_inside_limits & endZ_inside_limits)
\end{lstlisting}


The formal property is expressed in Computation Tree Logic (CTL) in Listing \ref{lst:CTLSpecNeg}.

\begin{lstlisting}[language=NuSMV, 
    caption={Reachability condition in CTL specification}, 
    label={lst:CTLSpecNeg}]
CTLSPEC 
AG !(endX_inside_limits & endY_inside_limits & endZ_inside_limits)
\end{lstlisting}

The specification:
\[
AG \neg ( endX\_inside\_limits \land endY\_inside\_limits \land endZ\_inside\_limits)
\]
is interpreted as: “On all computation paths (A), in all states globally 
(G), it is never the case that 
the end-effector simultaneously satisfies all three coordinate 
bounds.” In other words, if the  target point defined as in  Listing \ref{lst:CheckPreparation},
the model checker proves that the robot can never reach the point
(10,15,25) within the given tolerance.

Conversely, if the property is violated, NuSMV produces a counterexample trace that demonstrates 
a specific sequence of joint configurations leading the end
-effector into the defined target region. This 
counterexample can be directly interpreted as a feasible motion 
plan achieving the desired reachability.

The implementation  of this process centers on the Counterexample class in python
and a set of utilities for filtering variables, reconstructing geometry, 
and plotting results. This class  parse NuSMV counterexample, 
extracts the variables of interest (end-effector coordinates,
per-block coordinates, and joint angles), and provides means to render the 
counterexample trace in 3D graph. It also exposes numerically clean sequences that can 
hand to a controller. 

Counterexample class shown in Appendix \ref{ch:counterexamplePy} 
it takes raw NuSMV output and exposes methods 
to transform it into structured sequences for visualization and control. 
The source file is short—easy to audit and adapt—but covers the full 
parsing loop: reading states, carrying forward unchanged values, projecting 
specific variables, computing angles in degrees, reconstructing robot 
geometry step-by-step, and plotting.

Parsing the counterexample text presents a challenge because NuSMV prints 
only the variables that change between consecutive states. 
To reconstruct complete state vectors, the \texttt{filter\_variables$()$} function
employs a carry-forward mechanism. 
It begins with a list of variable names of interest, such as the coordinates of 
each block or the end-effector, and then iterates over the textual representation 
of each NuSMV state. When a variable assignment is found, its value is 
updated in the active state vector; otherwise, the previous value is 
retained. This mechanism ensures that every time step in the resulting sequence 
contains a full snapshot of all relevant variables, even those that remained 
constant. Once the entire sequence is reconstructed, the program converts 
the data into numerical form, yielding a well-defined matrix of dimension 
steps $\times$ variables suitable for quantitative analysis.
The resulting data provide a clear temporal evolution of the robot configuration.

When angular values are needed, the same mechanism is applied to 
orientation variables in function \texttt{get\_angles$()$}, 
which are scaled by a discretization factor to 
convert symbolic or integer representations into physical degrees. 
This transformation connects the abstract state space of the model checker 
to the metric quantities used in real robot kinematics.

Beyond parsing, the class includes visualization function \texttt{plot\_trace$()$} that convert 
the numeric data into three-dimensional figures. These plots display the 
complete end-effector path as well as the robot’s configuration at 
selected steps, an example of the plot is shown in Figure \ref{fig:3dCounterexampleTrace}.
Each step is represented as a connected chain of block 
positions, preserving the spatial structure of the model. The visualization enforces 
uniform scaling across all axes to maintain geometric fidelity, and 
the plotted points are labeled to correspond with discrete time indices. 

\begin{figure}[h!]
    \centering
    \includegraphics[height=12cm]{3dCounterexampleTrace.png} 
    \caption{Example trace (blue) plot for target point coordinates (10,13,30) 
    and last trace robot position (orange)}
    \label{fig:3dCounterexampleTrace}
\end{figure}

All paths and parameter values, such as the number of robot modules or 
the NuSMV executable location, are configurable. The implementation follows 
a transparent design that favors explicitness over abstraction, which makes it accessible to 
researchers who wish to adapt it to new model structures or naming conventions.

Once the counterexample has been parsed and verified, the data can be 
integrated into a robotic control pipeline. 

Data from counterexample can be used directly as open-loop trajectories, in which 
each time step of the counterexample corresponds to a discrete command for 
the robot’s end-effector or joint configuration. 
They may also serve as reference paths for closed-loop control (with feedback from sensors),
allowing the robot to reproduce formally verified behaviors under physical constraints. 

A typical workflow begins with 
the definition of the robot structure and target point in a NuSMV model 
file. The property under verification expresses the negation of the 
desired reachability condition -- formally, the statement that the robot 
can never reach the point. If NuSMV finds that this negated property is 
false, it outputs a counterexample that shows precisely how the robot 
can, in fact, reach the point. The Python parser then extracts from this 
output the sequence of coordinates or angles, which can be sent to 
the robot controller as a reference trajectory. 
Visualization confirms that the motion reaches the target, and further analysis can evaluate the 
dynamic parameters, adjust path smoothness, or joint feasibility of the resulting motion.

By systematically reconstructing state trajectories from NuSMV output, 
visualizing them, and exposing them for control, the Counterexample class enables the 
practical reuse of verification results in robotic system experiments.
Its structure is  extensible for larger models, alternative naming conventions, 
or  integration into simulation. 


% \subsection{Evaluation}
% TODO: plan experiments: 
% 1. Test with predefined assembly process (search, move, repeat) 
% with random distribution of the parts\\
% 2. Test with predefined assembly process (...) with, criteria and analysis 
% Criteria: success rate, average time, disagreement with control alg\\
%
% Experiment 1:\\ 
% Subject: local assembly with shortest path control algorighm\\
% Setup: random distribution of blocks and obstacles in 3D (for 0-15 obstacles) 
% Measure values: average assembly time, average success rate
%
%
% Experiment 2:\\ 
% Subject: local assembly with path control algorighm and runtime verification \\
% Setup: random distribution of blocks and obstacles in 3D (for 0-15 obstacles) 
% Measure values: average assembly time, average success rate
%
%
% Experiment 3:\\ 
% Subject: local assembly with runtime verification and counterexample based control algorighm\\
% Setup: random distribution of blocks and obstacles in 3D (for 0-15 obstacles) 
% Measure values: average assembly time, average success rate


\newpage




